\documentclass[pdflatex,sn-mathphys,Numbered]{sn-jnl}

\usepackage{graphicx}
\usepackage{amsmath,amssymb,amsfonts}
\usepackage{amsthm}
\usepackage{mathrsfs}
\usepackage[title]{appendix}
\usepackage{xcolor}
\usepackage{textcomp}
\usepackage{manyfoot}
\usepackage{booktabs}

\theoremstyle{thmstyleone}
\newtheorem{theorem}{Theorem}[section]
\newtheorem{proposition}[theorem]{Proposition}
\newtheorem{lemma}[theorem]{Lemma}
\newtheorem{corollary}[theorem]{Corollary}
\newtheorem{question}[theorem]{Question}
\theoremstyle{thmstyletwo}
\newtheorem{remark}[theorem]{Remark}
\newtheorem{definition}[theorem]{Definition}
\newtheorem{convention}[theorem]{Convention}

\newcommand{\MH}{\mathrm{MH}}
\newcommand{\PH}{\mathrm{PH}}
\newcommand{\RH}{\mathrm{RH}}
\newcommand{\MC}{\mathrm{MC}}
\newcommand{\QQ}{\mathbb{Q}}
\newcommand{\ZZ}{\mathbb{Z}}
\newcommand{\dwl}{\mathrm{DWL}}
\newcommand{\dwlp}{\mathrm{DWL}^{+}}

% AUTO-GENERATED by src/gen_tables.py -- do not edit by hand

\newcommand{\tabclasscounts}{%
\begin{tabular}{lrrrrrrr}
\toprule
$n$ & digraphs & D-WL & $\MH_{k\le 3}$ & dist & $\PH_{\le 2}$ & $\RH_{\le 3}$ & all three \\
\midrule
3 & 16 & 16 & 11 & 11 & 4 & 3 & 11 \\
4 & 218 & 215 & 74 & 48 & 6 & 5 & 74 \\
5 & 9608 & 9567 & 1415 & 237 & 11 & 8 & 1417 \\
\bottomrule
\end{tabular}}

\newcommand{\tabsplits}{%
\begin{tabular}{lrrr}
\toprule
relation & $n=3$ & $n=4$ & $n=5$ \\
\midrule
D-WL-equal, $\MH$-different & 0 & 2 & 25 \\
D-WL-equal, $\PH$-different & 0 & 1 & 6 \\
D-WL-equal, dist-different & 0 & 2 & 22 \\
dist-equal, D-WL-different & 4 & 33 & 210 \\
$\MH$-equal, D-WL-different & 3 & 38 & 1197 \\
$\MH$-equal, $\PH$-different & 0 & 0 & 2 \\
$\PH$-equal, $\MH$-different & 2 & 4 & 8 \\
$\PH$-equal, $\RH$-different & 0 & 1 & 3 \\
$\MH{+}\PH{+}\RH$-equal, D-WL-different & 3 & 38 & 1197 \\
$\MH{+}\PH$-equal, $\RH$-different & 0 & 0 & 0 \\
\bottomrule
\end{tabular}}

\newcommand{\mhcfour}{$\MH_{0,0}=4$, $\MH_{1,1}=4$, $\MH_{2,4}=4$, $\MH_{3,5}=4$}

\newcommand{\mhdigons}{$\MH_{0,0}=4$, $\MH_{1,1}=4$, $\MH_{2,2}=4$, $\MH_{3,3}=4$}

\newcommand{\mhcirc}{$\MH_{0,0}=5$, $\MH_{1,1}=10$, $\MH_{2,2}=10$, $\MH_{2,3}=10$, $\MH_{3,3}=10$, $\MH_{3,4}=30$}

\newcommand{\mhcirctwin}{$\MH_{0,0}=5$, $\MH_{1,1}=10$, $\MH_{2,2}=10$, $\MH_{2,3}=6$, $\MH_{3,3}=6$, $\MH_{3,4}=26$}

\newcommand{\beyonddistpairs}{\texttt{\&DKMME?} vs \texttt{\&DKLMI?}: arcs $\{02, 03, 13, 14, 20, 24, 30, 31, 41, 42\}$ vs $\{02, 03, 13, 14, 21, 24, 30, 31, 40, 42\}$ \par \texttt{\&DMMMU?} vs \texttt{\&DKVBZ?}: arcs $\{02, 03, 04, 13, 14, 20, 24, 30, 31, 34, 41, 42\}$ vs $\{02, 03, 12, 14, 20, 21, 31, 32, 34, 40, 42, 43\}$ \par \texttt{\&DKMMN?} vs \texttt{\&DKVYX?}: arcs $\{02, 03, 13, 14, 20, 24, 30, 31, 40, 41, 42, 43\}$ vs $\{02, 03, 12, 14, 20, 21, 23, 24, 31, 34, 40, 43\}$}


\raggedbottom

\begin{document}

\title[Directed Weisfeiler--Leman refinement and digraph homology]{Directed
Weisfeiler--Leman refinement versus the invariants of the magnitude--path
spectral sequence}

\author*[1]{\fnm{Krishna} \sur{Harish}}\email{krishnaharish2009@gmail.com}

\affil*[1]{\orgname{Elkins High School}, \orgaddress{\city{Missouri City},
\state{TX}, \country{USA}}}

\abstract{Magnitude homology, path homology, and reachability homology are
organized by the magnitude--path spectral sequence (MPSS) of a digraph:
magnitude homology is its first page, path homology sits on an axis of its
second page, and reachability homology is its target. Directed
Weisfeiler--Leman color refinement (D-WL) is the standard upper bound for
the separating power of message-passing neural networks on digraphs. We
give the first systematic comparison of the two families, the directed
counterpart of the known placement of persistent homology in the undirected
WL hierarchy. First, the families are incomparable: the directed $4$-cycle
and the disjoint union of two digons are D-WL-equivalent but are separated
by all three homology theories, while an in-star and a directed path of
length two are separated by D-WL although all three homology theories
agree; by exhaustive exact computation over all $9{,}842$ digraphs on
three, four, and five vertices, both witnesses are minimal. The incomparability
persists for the reciprocity-aware refinement $\dwlp$, with minimal witness
exactly at six vertices. Second, the bigraded rank function of the first
page of the MPSS does not determine the second page: any two digraphs in
which every vertex is a source or a sink and which share vertex and arc
counts have isomorphic magnitude homology in every bidegree, yet path
homology recovers the Betti numbers of the underlying graph; on at most
four vertices, exhaustively, no such rank-equal pair with distinct path
homology exists, so five vertices is minimal. Third, magnitude homology
separates D-WL-equivalent digraphs with equal directed-distance multisets,
so its excess power over refinement is not distance bookkeeping; the
minimal witnesses (five vertices, none at four or fewer) pit the symmetric
digraph of the $5$-cycle against non-circulant partial orientations.
All homology ranks are computed exactly over $\QQ$; we release the
accompanying toolkit, to our knowledge the first for magnitude, path, and
reachability homology of general digraphs in a single exact-arithmetic
framework.}

\keywords{magnitude homology, path homology, reachability homology,
Weisfeiler--Leman, digraphs, graph neural networks, expressivity}

\maketitle

\section{Introduction}\label{sec:intro}

Two research communities have, largely independently, built graph-level
invariants of directed graphs.

On the applied-topology side, a family of homology theories for digraphs
has crystallized over the past decade: the path homology of Grigor'yan,
Lin, Muranov, and Yau \cite{GLMY12,GLMY14}, sensitive to the composability
of arcs; the magnitude homology of Hepworth--Willerton and
Leinster--Shulman \cite{HepworthWillerton17,LeinsterShulman21},
categorifying Leinster's magnitude \cite{Leinster13} and defined for any
digraph through its shortest-path quasimetric; and the reachability
homology of Hepworth and Roff \cite{HepworthRoffIMRN}. These are not
independent inventions: Asao \cite{Asao23} constructed a spectral sequence,
the magnitude--path spectral sequence (MPSS), whose first page is magnitude
homology and whose second page contains path homology, and Hepworth and
Roff \cite{HepworthRoff24,HepworthRoffIMRN} showed that it converges to
reachability homology. The MPSS thus organizes the three theories into
pages of a single object.

On the machine-learning side, the separating power of message-passing
graph neural networks is bounded by Weisfeiler--Leman color refinement
\cite{Xu19,Morris19}; for digraphs the relevant procedure is directed
color refinement (here D-WL), which refines vertex colors by the multisets
of in-neighbor and out-neighbor colors separately \cite{GroheColRef,
Jaume19}. For undirected graphs, the relationship between persistent
homology and the WL hierarchy has been worked out in detail by Ballester
and Rieck \cite{BallesterRieck25}, but that study is explicitly restricted
to undirected graphs, and we are not aware of any placement of the
digraph homology theories above relative to directed refinement. The
closest prior results are single-invariant comparisons within topology,
such as the digraph pair of Chaplin, Harrington, and Tillmann
\cite{ChaplinHT24} distinguishing path homology from directed flag
complex homology.

This paper carries out the comparison. Our contributions are as follows;
throughout, $\MH$, $\PH$, $\RH$ denote magnitude, path, and reachability
homology over $\QQ$ (Section~\ref{sec:background}), and all computational
claims are exact rational computations, exhaustive over all digraphs on
three, four, and five vertices up to isomorphism ($16 + 218 + 9{,}608 =
9{,}842$ digraphs; the four digraphs on at most two vertices are pairwise
D-WL-distinguished, so all minimality thresholds are meaningful).

\begin{enumerate}
\item \textbf{Incomparability} (Theorem~\ref{thm:incomparable}). D-WL and
the MPSS invariants each separate digraph pairs the other cannot. The
directed $4$-cycle versus two digons is D-WL-equivalent but separated by
each of $\MH$, $\PH$, $\RH$; an in-star versus a directed $2$-path is
D-WL-separated while all three homology theories agree. Both witnesses are
minimal by exhaustive search. For the reciprocity-aware refinement
$\dwlp$, which additionally distinguishes reciprocal neighbors, the
homology-beats-refinement witness moves to the directed $6$-cycle versus
two directed triangles, minimal at exactly six vertices
(Theorem~\ref{thm:dwlplus}).
\item \textbf{The first page does not determine the second}
(Theorem~\ref{thm:e1e2}). On the family of digraphs in which every vertex
is a source or a sink, all magnitude homology is concentrated in bidegrees
$(0,0)$ and $(1,1)$ with ranks $(n, m)$, while path homology and
reachability homology compute the Betti numbers of the underlying
undirected graph. Consequently the bigraded ranks of the $E^1$ page of the
MPSS do not determine its $E^2$ page, its limit, or even weak
connectivity. Exhaustively, no pair of digraphs on at most four vertices
with equal $\MH$ ranks has distinct path homology, so the minimal such
pair, the in-star $K_{4\to 1}$ versus $K_{2\to 2}\sqcup\ast$, lives on
five vertices.
\item \textbf{Beyond distances} (Theorem~\ref{thm:beyonddist}). Magnitude
homology ranks do not determine the directed-distance multiset
(Proposition~\ref{prop:mhnotdist}), and conversely magnitude homology
separates D-WL-equivalent pairs with \emph{equal} distance multisets:
minimally, on five vertices, the symmetric digraph of the $5$-cycle
against a partial orientation with three digons, where path homology even
distinguishes a $1$-dimensional from a $2$-dimensional hole. All such
witnesses on at most five vertices carry unequal digon counts, and the
reciprocity-aware $\dwlp$ together with distances subsumes the homology
package in this range; whether magnitude homology exceeds
$(\dwlp,\text{distances})$ on larger digraphs is open
(Question~\ref{q:dwlplusdist}).
\item \textbf{The landscape} (Section~\ref{sec:landscape}). We tabulate
the partition of all $9{,}842$ digraphs on at most five vertices induced
by each invariant and each pairwise refinement relation, including that
$\PH$ does not determine $\RH$, while $\MH$ and $\PH$ jointly determine
$\RH$ throughout this range (Question~\ref{q:mhphrh}).
\item \textbf{Software.} The exact-arithmetic toolkit and all data are
released; to our knowledge it is the first public implementation
computing magnitude homology of general digraphs, and the first computing
the three MPSS-related theories in one framework. (For undirected graphs,
computations of graph magnitude homology appear in the literature, for
example via algebraic Morse theory \cite{Gu18}.)
\end{enumerate}

Beyond the specific statements, the results calibrate expectations for
homology-based features in directed graph learning: the MPSS invariants
are neither refinements nor coarsenings of the D-WL bound that governs
standard directed message passing, so they are candidates for genuinely
complementary features, but they also miss elementary refinement
information (weak connectivity is invisible to magnitude homology, while
degree profiles separate what all three theories jointly cannot). We
return to this in Section~\ref{sec:discussion}.

\subsection*{Relation to prior work}

Ballester and Rieck \cite{BallesterRieck25} place persistent homology of
filtered clique complexes within the (undirected) WL hierarchy; our
results are the directed analogue at the level of unfiltered digraph
homology theories, and Section~\ref{sec:discussion} discusses the
persistent directed extension as future work. Path-based neural
architectures \cite{Michel23} exceed 1-WL by aggregating over paths but
involve no homology. Within topology, comparisons among directed
invariants exist \cite{ChaplinHT24,CaputiMR25}, but none against
refinement hierarchies. On the computational side, Menara and Manzoni
\cite{MenaraManzoni24} analyze the parameterized complexity of Eulerian
magnitude homology of graphs; our sweep is feasible precisely because the
digraphs are small and the arithmetic is exact. The MPSS papers
\cite{Asao23,HepworthRoff24,HepworthRoffIMRN} study the algebraic
structure of the tower; our Theorem~\ref{thm:e1e2} is, to our knowledge,
the first explicit witness that the rank function of $E^1$ does not
determine $E^2$, complementing their structural results. Persistent
magnitude homology has just been given foundations \cite{BiFengLiWu26};
its separating power is untouched and inherits our lower bounds.

\section{Background and conventions}\label{sec:background}

\begin{convention}
A \emph{digraph} $G = (V, A)$ is a finite set $V$ with a set
$A \subseteq V \times V$ of ordered pairs (\emph{arcs}) with no loops;
$n = |V|$, $m = |A|$. A \emph{digon} is a pair of reciprocal arcs
$\{(u,v),(v,u)\}$. The \emph{quasimetric} $d(x,y) \in \{0,1,2,\dots\}
\cup \{\infty\}$ is the directed shortest-path distance. All homology is
taken with coefficients in $\QQ$ and we record ranks. Digraph
enumeration is up to isomorphism via nauty \cite{nauty}; the counts
$16$, $218$, $9{,}608$ for $n = 3,4,5$ match OEIS A000273.
\end{convention}

\subsection{Directed color refinement}

D-WL iteratively refines a vertex coloring: all vertices start with a
common color, and the color of $v$ at round $t{+}1$ is the (injectively
encoded) triple of its color at round $t$, the multiset of round-$t$
colors of its in-neighbors, and the multiset of round-$t$ colors of its
out-neighbors \cite{GroheColRef,Jaume19}. Since a partition of an $n$-set
can strictly refine at most $n-1$ times, the partition stabilizes within
$n$ rounds. Two digraphs are \emph{D-WL-equivalent} if the color
histograms agree at every round (colors encoded canonically across
graphs). This is the separating power that bounds standard directed
message-passing networks, in the same way that 1-WL bounds GNNs on
undirected graphs \cite{Xu19,Morris19}.

We will also use a natural strengthening. The \emph{reciprocity-aware}
refinement $\dwlp$ refines by the triple of multisets over in-only
neighbors, out-only neighbors, and reciprocal neighbors (those joined to
$v$ by a digon). $\dwlp$ refines D-WL; it corresponds to message passing
that treats reciprocated and unreciprocated arcs as different edge types,
as many practical digraph architectures do.

\subsection{Magnitude homology}

Following \cite{HepworthWillerton17,LeinsterShulman21}, applied to the
quasimetric space $(V, d)$: for $k \ge 0$ and $l \ge 0$, the chain group
$\MC_{k,l}(G)$ is the $\QQ$-span of tuples $(x_0, \dots, x_k)$ of
vertices with $x_{i-1} \ne x_i$, $d(x_{i-1}, x_i) < \infty$, and
$\sum_{i=1}^{k} d(x_{i-1}, x_i) = l$. The differential is
$\partial = \sum_{i=1}^{k-1} (-1)^i \partial_i$, where $\partial_i$
deletes $x_i$ if $d(x_{i-1},x_i) + d(x_i,x_{i+1}) = d(x_{i-1},x_{i+1})$
and is zero on the tuple otherwise. Magnitude homology $\MH_{k,l}(G)$ is
the homology of the $l$-graded piece. Its graded Euler characteristic
recovers the magnitude power series; we use this identity as an
end-to-end correctness test of our implementation
(Section~\ref{sec:landscape}). By $\MH$ \emph{ranks} we always mean the
function $(k,l) \mapsto \dim_\QQ \MH_{k,l}$.

\subsection{Path homology}

Following \cite{GLMY12,GLMY14}: an \emph{allowed} $p$-path is a vertex
sequence $(x_0,\dots,x_p)$ each consecutive pair of which is an arc. On
the span of \emph{regular} paths (consecutive entries distinct) one takes
the alternating face boundary with irregular faces set to zero, and
$\Omega_p = \{\omega \in \operatorname{span}(A_p) : \partial\omega \in
\operatorname{span}(A_{p-1})\}$ is the space of $\partial$-invariant
allowed chains. Path homology $\PH_p(G)$ is the homology of
$(\Omega_\ast, \partial)$. We compute $\PH_0, \PH_1, \PH_2$.

\subsection{Reachability homology and the MPSS}

Reachability homology $\RH_\ast(G)$ \cite{HepworthRoffIMRN} is the
homology of the nerve of the reachability preorder ($x \preceq y$ iff
$d(x,y) < \infty$). The preorder is equivalent, as a category, to its
condensation poset (strong components ordered by reachability), so
$\RH_\ast(G)$ is the simplicial homology of the order complex of the
condensation; this is how we compute it, in dimensions $0$ to $3$.
Asao \cite{Asao23} constructs the MPSS with $E^1 \cong \MH$; the diagonal
of its $E^2$ page is reduced path homology, and the sequence converges to
reachability homology \cite{HepworthRoff24,HepworthRoffIMRN}. We only use
this organizational fact: $\MH$, $\PH$, $\RH$ are, respectively, the
first page, part of the second page, and the target of one spectral
sequence.

\subsection{Windows, and what is absolute}\label{subsec:windows}

Sweep-level computations use the windows $k \le 3$ (all $l$) for $\MH$,
$p \le 2$ for $\PH$, and dimensions $\le 3$ for $\RH$. An
\emph{inequality} of invariants observed in a window is absolute: the
full invariants differ. An \emph{equality} claim is, in general, relative
to the window and is stated as such; the star witnesses of
Theorem~\ref{thm:e1e2} are the exception, since for them the full
magnitude homology is determined by hand in all bidegrees.

\section{Structural lemmas}\label{sec:lemmas}

\begin{lemma}\label{lem:regular}
Let $G$ and $H$ be digraphs on $n$ vertices in which every vertex has
out-degree $a$ and in-degree $b$. Then $G$ and $H$ are D-WL-equivalent.
If moreover neither contains a digon, then $G$ and $H$ are
$\dwlp$-equivalent.
\end{lemma}

\begin{proof}
By induction on the round: all vertices of both digraphs share the
initial color, and if all share a color $c_t$ at round $t$, then every
vertex of either digraph has the same signature
$(c_t, \{c_t^{(b)}\}, \{c_t^{(a)}\})$, hence the same round-$(t{+}1)$
color. The histograms therefore agree at every round. In the digon-free
case the reciprocal multiset is empty for every vertex and the same
induction applies to $\dwlp$.
\end{proof}

\begin{lemma}\label{lem:mh1}
For every digraph $G$: $\MH_{0,0}(G) \cong \QQ^{n}$,
$\MH_{1,1}(G) \cong \QQ^{m}$, and $\MH_{1,l}(G) = 0$ for all $l \ge 2$.
\end{lemma}

\begin{proof}
Degree $0$ is immediate. $\MC_{1,l}$ is spanned by ordered pairs at
distance exactly $l$, and the outgoing differential vanishes on
$\MC_{1,\ast}$ (there are no interior vertices to delete). At $l = 1$
the pairs at distance $1$ are exactly the arcs and $\MC_{2,1} = 0$, so
$\MH_{1,1} \cong \QQ^m$. For $l \ge 2$, let $d(x,y) = l$ and let $z$ be
the first interior vertex of a geodesic from $x$ to $y$; then
$d(x,z) = 1$ and $d(z,y) = l - 1$, so $(x,z,y) \in \MC_{2,l}$ and
$\partial(x,z,y) = -(x,y)$. Hence the incoming differential is
surjective and $\MH_{1,l} = 0$.
\end{proof}

Thus all information carried by magnitude homology beyond the vertex and
arc counts lives in homological degrees $k \ge 2$. The next observation
identifies a family where those degrees vanish identically.

\begin{lemma}\label{lem:sourcesink}
Call a digraph \emph{layered} if every vertex is a source or a sink
(no vertex has both an incoming and an outgoing arc). If $G$ is layered
with $n$ vertices, $m$ arcs, and $c$ weakly connected components, then:
\begin{enumerate}
\item $\MC_{k,l}(G) = 0$ for all $k \ge 2$, hence $\MH(G)$ has rank $n$
at $(0,0)$, rank $m$ at $(1,1)$, and rank $0$ elsewhere;
\item $\PH_0(G) \cong \QQ^{c}$, $\PH_1(G) \cong \QQ^{\,m - n + c}$, and
$\PH_p(G) = 0$ for $p \ge 2$;
\item $\RH_0(G) \cong \QQ^{c}$, $\RH_1(G) \cong \QQ^{\,m - n + c}$, and
$\RH_p(G) = 0$ for $p \ge 2$.
\end{enumerate}
\end{lemma}

\begin{proof}
(1) A tuple $(x_0, x_1, x_2, \dots)$ in $\MC_{k,l}$ with $k \ge 2$
requires $d(x_0, x_1) < \infty$ and $d(x_1, x_2) < \infty$ with
$x_0 \ne x_1 \ne x_2$, so $x_1$ has both an incoming and an outgoing arc
(any finite positive distance requires at least one arc out of the source
and into the target), contradicting layeredness.
(2) In a layered digraph there are no allowed $2$-paths, so
$\Omega_p = 0$ for $p \ge 2$, $\Omega_1$ is spanned by the arcs, and
$\Omega_0$ by the vertices. The complex
$0 \to \QQ^m \xrightarrow{\;\partial\;} \QQ^n \to 0$ is the simplicial
chain complex of the underlying undirected graph, whose homology is
$(\QQ^c, \QQ^{m-n+c})$.
(3) Layeredness forbids digons and directed $2$-paths, so the
reachability preorder is the height-one poset whose comparabilities are
the arcs, all strong components are singletons, and the order complex is
the underlying undirected graph (its edges are the arcs; there are no
chains of length two). The claim follows as in (2).
\end{proof}

\begin{proposition}\label{prop:mhnotdist}
Magnitude homology ranks do not determine the multiset of directed
distances. Minimal witnesses exist on three vertices: the in-star
$S$ ($u \to w \leftarrow v$) and the directed path
$P$ ($u \to w \to v$) satisfy $\MH_{k,l}(S) \cong \MH_{k,l}(P)$ for all
$(k,l)$, but the finite distance multisets are $\{1,1\}$ and
$\{1,1,2\}$.
\end{proposition}

\begin{proof}
$S$ is layered, so by Lemma~\ref{lem:sourcesink} its $\MH$ is
$(3, 2)$ at bidegrees $(0,0), (1,1)$ and zero elsewhere. In $P$, the only
tuples beyond degree one are built from the composable pair
$u \to w \to v$: $\MC_{2,2}(P)$ is spanned by $(u,w,v)$, and
$\MC_{k,l}(P) = 0$ for $k \ge 3$ (nothing leaves $v$). Since
$d(u,w) + d(w,v) = 2 = d(u,v)$, we get
$\partial(u,w,v) = -(u,v) \ne 0$, so the single degree-$2$ chain is not a
cycle and it kills the single class of $\MC_{1,2}$: explicitly
$\MH_{1,2}(P) = 0$, $\MH_{2,2}(P) = 0$. Hence $\MH(P)$ also equals
$(3,2)$ at $(0,0), (1,1)$ and zero elsewhere, while $d(u,v) = 2$ appears
only in $P$.
\end{proof}

\section{Incomparability}\label{sec:incomparable}

\begin{theorem}\label{thm:incomparable}
D-WL and the family $\{\MH, \PH, \RH\}$ are incomparable in separating
power:
\begin{enumerate}
\item[(a)] The directed $4$-cycle $C_4$ and the disjoint union of two
digons $D_2 \sqcup D_2$ are D-WL-equivalent, but each of $\MH$, $\PH$,
$\RH$ separates them. No such pair exists on at most three vertices.
\item[(b)] The in-star $S$ and directed path $P$ of
Proposition~\ref{prop:mhnotdist} satisfy
$\MH(S) \cong \MH(P)$ (all bidegrees), $\PH_\ast(S) \cong \PH_\ast(P)$,
$\RH_\ast(S) \cong \RH_\ast(P)$, but D-WL separates them at round one.
\end{enumerate}
\end{theorem}

\begin{proof}
(a) Both digraphs are $1$-in-$1$-out regular on four vertices, so
Lemma~\ref{lem:regular} gives D-WL-equivalence. $\RH$: $C_4$ is strongly
connected, so its reachability preorder is equivalent to a point and
$\RH_0 = \QQ$; the union of digons has two strong components, mutually
unreachable, giving $\RH_0 = \QQ^2$. $\PH$: $\PH_0$ counts weak
components ($1$ versus $2$); moreover $\PH_1(C_4) = \QQ$ (the four arcs
support the cycle class $\sum e_i$, and $\Omega_2 = 0$ since every
allowed $2$-path $(x, x{+}1, x{+}2)$ has the non-allowed face
$(x, x{+}2)$, and these four faces are distinct, so no linear combination
cancels them), while $\PH_1(D_2 \sqcup D_2) = 0$: for a single digon on
$\{x,y\}$, $\Omega_1$ is spanned by $e_{xy}, e_{yx}$ with $\partial$ of
rank $1$, and the surviving cycle $e_{xy} + e_{yx}$ is a boundary, since
the allowed $2$-path $(x,y,x)$ is $\partial$-invariant with
$\partial(x,y,x) = e_{yx} + e_{xy}$ (the middle face $(x,x)$ is
irregular, hence zero), so $\PH_1 = 0$ for each digon and their disjoint
union.
$\MH$: the machine computation (exact over $\QQ$, Table~\ref{tab:star})
gives $\MH_{2,4}(C_4) = \QQ^4$ while $\MH_{2,4}(D_2 \sqcup D_2) = 0$;
alternatively, the digon union is the symmetric digraph of a forest-like
graph ($2 K_2$), whose magnitude homology is diagonal
\cite{HepworthWillerton17}, while $C_4$ has nonzero off-diagonal ranks.
Since a rank inequality observed in any bidegree is absolute
(Section~\ref{subsec:windows}), all three separations hold. Minimality is
the exhaustive $n = 3$ computation (Table~\ref{tab:splits}: zero D-WL
classes split by any homology at $n = 3$).

(b) Magnitude homology equality is Proposition~\ref{prop:mhnotdist}.
$\PH$: both digraphs are weakly connected with two arcs whose boundaries
are linearly independent, and in both $\Omega_2 = 0$ ($S$ has no allowed
$2$-path; in $P$ the single allowed $2$-path $(u,w,v)$ has non-allowed
face $(u,v)$), so $\PH_\ast = (\QQ, 0, 0)$ for both. $\RH$: the
reachability posets are $u, v \prec w$ (order complex: a path, hence
contractible) and $u \prec w \prec v$ with $u \prec v$ (order complex: a
$2$-simplex, contractible), so $\RH_\ast = (\QQ, 0, 0, 0)$ for both.
D-WL separates at round one: the in-degree multisets are
$\{0,0,2\}$ and $\{0,1,1\}$.
\end{proof}

\begin{table}[t]
\centering
\caption{Magnitude homology ranks of the star witnesses of
Theorem~\ref{thm:incomparable}(a) and Theorem~\ref{thm:beyonddist}
(machine-generated; nonzero bidegrees $(k,l)$ with $k \le 3$).}
\label{tab:star}
\begin{tabular}{ll}
\toprule
digraph & nonzero $\MH_{k,l}$, $k \le 3$ \\
\midrule
directed $C_4$ & \mhcfour \\
$D_2 \sqcup D_2$ & \mhdigons \\
symmetric $C_5$ (\texttt{\&DKMME?}) & \mhcirc \\
twin (\texttt{\&DKLMI?}) & \mhcirctwin \\
\bottomrule
\end{tabular}
\end{table}

The refinement $\dwlp$ sees digons, and the witness of part (a) has
digons on one side only. The incomparability nevertheless persists, one
size up, and exactly there:

\begin{theorem}\label{thm:dwlplus}
The directed $6$-cycle $C_6$ and the disjoint union $C_3 \sqcup C_3$ of
two directed triangles are $\dwlp$-equivalent, and each of $\MH$, $\PH$,
$\RH$ separates them. Exhaustively, no $\dwlp$-equivalent pair of
digraphs on at most five vertices is separated by $\MH$ (ranks
$k \le 3$) or by $\PH_{\le 2}$; hence six vertices is minimal.
\end{theorem}

\begin{proof}
Both digraphs are digon-free and $1$-in-$1$-out regular on six vertices,
so Lemma~\ref{lem:regular} gives $\dwlp$-equivalence. Separations:
$\PH_0$ and $\RH_0$ count weak components and strong condensation
components ($1$ versus $2$ each); $\PH_\ast(C_6) = (\QQ, \QQ, 0)$ and
$\PH_\ast(C_3 \sqcup C_3) = (\QQ^2, \QQ^2, 0)$ (directed cycles have
$\PH_1 = \QQ$ \cite{GLMY12}; verified exactly by our implementation);
$\MH_{1,\ast}$ agrees (Lemma~\ref{lem:mh1}) but for instance
$\MC_{1,3}$-level distances differ and the machine gives, absolutely,
$\MH_{2,3}(C_3 \sqcup C_3) = \QQ^{6} \ne \QQ^{0} = \MH_{2,3}(C_6)$.
The exhaustive statement is a direct computation over all $9{,}842$
digraphs on $3 \le n \le 5$, together with the trivial check that the
four digraphs on $n \le 2$ are pairwise D-WL-distinguished
(Section~\ref{sec:landscape}).
\end{proof}

\section{The first page does not determine the second}\label{sec:e1e2}

\begin{theorem}\label{thm:e1e2}
Let $X = K_{4 \to 1}$ (four sources, one sink, all four arcs) and
$Y = K_{2\to 2} \sqcup \ast$ (complete bipartite orientation plus an
isolated vertex), both on five vertices with four arcs. Then:
\begin{enumerate}
\item[(a)] $\MH_{k,l}(X) \cong \MH_{k,l}(Y)$ for \emph{all} $(k,l)$;
\item[(b)] $\PH_\ast(X) = (\QQ, 0, 0) \ne (\QQ^2, \QQ, 0) = \PH_\ast(Y)$
and $\RH_\ast(X) = (\QQ, 0, 0, 0) \ne (\QQ^2, \QQ, 0, 0) =
\RH_\ast(Y)$.
\end{enumerate}
Hence the bigraded rank function of the first page of the MPSS determines
neither its second page nor its target, and magnitude homology does not
determine weak connectivity. Exhaustively, no pair of digraphs on at most
four vertices with equal $\MH$ ranks ($k \le 3$, all $l$) has distinct
$\PH_{\le 2}$, so five vertices is minimal.
\end{theorem}

\begin{proof}
Both digraphs are layered with $(n, m) = (5, 4)$;
Lemma~\ref{lem:sourcesink}(1) gives (a), and
Lemma~\ref{lem:sourcesink}(2),(3) give (b) with $(c, m{-}n{+}c)$ equal to
$(1, 0)$ for $X$ and $(2, 1)$ for $Y$. For the MPSS statement: the $E^1$
pages have equal bigraded ranks by (a), while the $E^2$ page contains
(reduced) path homology \cite{Asao23,HepworthRoff24} and the target is
reachability homology \cite{HepworthRoffIMRN}, both of which differ by
(b). Minimality is the exhaustive $n \le 4$ computation
(Table~\ref{tab:splits}); note that an equality of the coarser $k \le 3$
rank function is implied by equality of the full rank function, so the
exhaustive statement covers the full-rank hypothesis a fortiori.
\end{proof}

\begin{remark}
Lemma~\ref{lem:sourcesink} yields infinitely many such pairs: any two
layered digraphs with equal $(n, m)$ and different $c$ (equivalently,
different cycle rank $m - n + c$) work. On five vertices with four arcs
the pair $(X, Y)$ above is the one produced by the exhaustive search. The
phenomenon is a rank statement about $E^1$: as pages-with-differentials
the two $E^1$ terms are of course not isomorphic, and the theorem shows
the differential carries essential information, in the strongest sense
that it is not recoverable from ranks.
\end{remark}

\begin{remark}\label{rem:truncation}
The sweep also finds pairs equal at $k \le 3$ that split at $k = 4$
(for example \texttt{\&DKCJN?} versus \texttt{\&DKWMM?}, whose $\MH_{4,4}$
ranks are $4$ and $6$): window equality is genuinely weaker than full
equality, which is why Theorem~\ref{thm:e1e2} is stated with a
hand-proved full-rank witness.
\end{remark}

\section{Beyond distance multisets}\label{sec:beyonddist}

D-WL does not determine directed distances, and distances do not
determine D-WL (Table~\ref{tab:splits}, rows 3 and 4); the pair
$(C_4, D_2 \sqcup D_2)$ realizes the first failure. The sharper question
is whether the homology theories exceed D-WL \emph{only} through
distances. On at most four vertices the answer is yes: every D-WL class
split by $\MH$ or $\PH$ is already split by the distance multiset. On
five vertices the answer is no:

\begin{theorem}\label{thm:beyonddist}
There exist digraphs on five vertices, minimally, that are
D-WL-equivalent and have equal multisets of directed distances but are
separated by $\MH$ and by $\PH$. Witness (machine-verified, exact): the
symmetric digraph $\overline{C_5}$ of the $5$-cycle
(arcs $x \to x \pm 2 \bmod 5$; five digons) and its partial-orientation
twin \texttt{\&DKLMI?} (arcs
$\{02, 03, 13, 14, 21, 24, 30, 31, 40, 42\}$; three digons). Both are
$2$-in-$2$-out regular, strongly connected with all $20$ ordered
distances in $\{1, 2\}$ ($10$ each), and D-WL-equivalent by
Lemma~\ref{lem:regular}; but $\MH_{2,3}$ has ranks $10$ versus $6$,
$\MH_{3,4}$ has ranks $30$ versus $26$, and
$\PH_\ast = (\QQ, \QQ, 0)$ versus $(\QQ, 0, \QQ)$: the twin trades the
$1$-dimensional hole of $\overline{C_5}$ for a $2$-dimensional one.
Exactly three such D-WL-and-distance-equivalent split classes exist on
five vertices, listed in Appendix~\ref{app:witnesses}.
\end{theorem}

\begin{proof}
Regularity, distances, digon counts, and D-WL-equivalence are finite
checks stated in the theorem (the distance claim follows since each
vertex has two out-neighbors and reaches the remaining two vertices in
two steps; Lemma~\ref{lem:regular} applies). The homology inequalities
are exact machine computations, absolute by
Section~\ref{subsec:windows}; the verification script in the repository
recomputes all invariants of this pair, including the per-witness Euler
characteristic identity, in seconds. Minimality is the exhaustive
$n \le 4$ statement above.
\end{proof}

\begin{remark}
$\overline{C_5}$ carries the magnitude homology of the $5$-cycle in the
sense of Hepworth--Willerton (a symmetric digraph's quasimetric is the
graph metric). The twin is not a circulant: a circulant digraph
$\mathrm{Circ}(5; S)$ with $|S| = 2$ has either zero or five digons
(according as $S \cap (-S)$ is empty or all of $S$), never three.
\end{remark}

\begin{remark}\label{rem:digons}
All three witness classes of Theorem~\ref{thm:beyonddist} pair digraphs
with different digon counts ($5$ versus $3$ in each), so the
reciprocity-aware $\dwlp$ separates each of them; indeed, exhaustively on
$n \le 5$, no $(\dwlp, \text{distance})$-equivalent pair is separated by
$\MH$ ($k \le 3$) at all. What magnitude homology detects here beyond
D-WL and distances is, at minimum, reciprocity structure; whether it
detects more on larger digraphs is Question~\ref{q:dwlplusdist}.
\end{remark}

\section{The landscape on at most five vertices}\label{sec:landscape}

\begin{table}[t]
\centering
\caption{Number of equivalence classes induced on the isomorphism classes
of digraphs by each invariant (machine-generated). ``dist'' is the
multiset of ordered directed distances; ``all three'' is the joint
$(\MH_{k\le3}, \PH_{\le 2}, \RH_{\le 3})$ class.}
\label{tab:classes}
\tabclasscounts
\end{table}

\begin{table}[t]
\centering
\caption{Refinement failures: number of equivalence classes of the first
invariant that split into two or more classes of the second
(machine-generated). A zero row is an exhaustive determination statement
for that $n$ and those windows.}
\label{tab:splits}
\tabsplits
\end{table}

Tables~\ref{tab:classes} and \ref{tab:splits} summarize the exhaustive
sweep. Beyond the theorems above, three features are worth recording.

First, D-WL is nearly complete on this range ($9{,}567$ classes out of
$9{,}608$ at $n = 5$), consistent with the classical almost-sure
completeness of color refinement; the homology theories are far coarser
as graph-level summaries. Their value is therefore not raw separating
power but the incomparability established above: what they see, D-WL can
miss entirely.

Second, within the MPSS family: $\PH$ does not determine $\RH$ (one
class splits at $n = 4$, three at $n = 5$), but $\MH$ and $\PH$
\emph{jointly} determine $\RH$ throughout the range, in the computed
windows; internally, the joint $(\MH, \PH, \RH)$ partition at $n = 5$
has exactly two more classes than the $\MH$ partition, and both extra
splits are the $E^1$-versus-$E^2$ witnesses of Section~\ref{sec:e1e2}.

\begin{question}\label{q:mhphrh}
Do the bigraded ranks of $\MH$ together with $\PH_\ast$ determine
$\RH_\ast$ for all digraphs? (True for all digraphs on at most five
vertices in the computed windows.)
\end{question}

\begin{question}\label{q:dwlplusdist}
Is there a pair of $\dwlp$-equivalent digraphs with equal directed
distance multisets separated by magnitude homology? (None exist on at
most five vertices; by Theorem~\ref{thm:dwlplus} the analogous question
without the distance condition has answer yes, minimally at $n = 6$.)
\end{question}

\subsection*{Methodology and verification}

All digraphs were enumerated with nauty (\texttt{geng | directg})
\cite{nauty}; counts match OEIS A000273. All boundary matrices are
integer matrices and all ranks are computed by fraction-free Gaussian
elimination over $\QQ$; no floating point is used anywhere. The
implementation is validated by a test suite of $69$ checks against
independently known values: the graded Euler characteristic of $\MH$
against the magnitude power series (computed independently from the
similarity matrix) across all gradings $l \le 6$ on seven digraphs and
graphs; Hepworth--Willerton diagonality and ranks for complete graphs;
$\partial^2 = 0$ for the regularized path boundary on all $218$
digraphs with $n = 4$; the standard path homology examples of
\cite{GLMY12} (directed cycles, filled and unfilled triangles and
squares); and contractibility statements for reachability homology. The
star witnesses of each theorem were additionally re-verified in isolation
with extended windows ($k \le 5$, $l \le 10$ for magnitude homology) and
per-witness Euler characteristic checks.

\section{Discussion}\label{sec:discussion}

\paragraph{Implications for directed graph learning.}
The invariants of the MPSS are neither stronger nor weaker than the
refinement bound governing standard directed message passing: they
separate D-WL-equivalent digraphs already at four vertices (and even
$\dwlp$-equivalent ones at six), while missing distinctions (weak
connectivity, degree profiles) that one round of refinement sees. For
feature design the conclusion is symmetric: homological features are
candidates to add genuinely new signal to directed MPNNs, exactly as
persistent homology is for undirected GNNs \cite{BallesterRieck25}, but
they cannot replace refinement. The beyond-distance witnesses sharpen
this: part of what magnitude homology adds is reciprocity structure,
which standard D-WL provably misses and which practical architectures
handle by edge-typing; Question~\ref{q:dwlplusdist} asks whether
magnitude homology sees beyond even that.

\paragraph{Implications within the MPSS.}
Theorem~\ref{thm:e1e2} shows in the strongest rank-level sense that the
$E^1$ differential of the MPSS carries information: two digraphs can have
identical $E^1$ rank functions and different $E^2$ pages, indeed
different limits. The layered family of Lemma~\ref{lem:sourcesink} may be
of independent interest as the locus where $E^1$ degenerates to counting
while $E^2$ recovers the topology of the underlying graph.

\paragraph{Limitations.}
Exhaustive claims are limited to $n \le 5$ and to the stated windows
(equalities only; inequalities are absolute). We compare against D-WL and
$\dwlp$, not against the higher levels of a directed WL hierarchy; a
directed analogue of the $k$-WL placement of \cite{BallesterRieck25}, and
the extension from unfiltered digraph homology to persistent path
homology \cite{ChowdhuryMemoli18} and persistent magnitude homology
\cite{BiFengLiWu26} on weighted digraphs, are natural next steps. Our
lower bounds transfer: any invariant that specializes to $\MH$, $\PH$, or
$\RH$ at a fixed filtration value inherits the corresponding
separations.

\paragraph{Data and code availability.}
The full implementation (exact-arithmetic magnitude, path, and
reachability homology of digraphs; D-WL and $\dwlp$; enumeration and
analysis drivers), all sweep outputs, and scripts reproducing every
number and table in this paper are available at
\url{https://github.com/krishnatheaverage/mpss-dwl-expressivity}.

\begin{appendices}

\section{Witness data}\label{app:witnesses}

Digraphs are given in nauty digraph6 format together with arc lists on
vertex set $\{0, \dots, 4\}$; $uv$ denotes the arc $u \to v$. The three
D-WL-and-distance-equivalent classes on five vertices split by magnitude
homology (Theorem~\ref{thm:beyonddist}) are represented by:

\beyonddistpairs

\medskip
In each pair the first digraph has five digons and the second has three.
The first pair is the star witness of Theorem~\ref{thm:beyonddist}; in
the remaining two pairs the path homology dimensions $(\PH_0, \PH_1,
\PH_2)$ agree, so the separation there is by magnitude homology alone.

\end{appendices}

\bibliography{refs}

\end{document}
